% Mathematical bridge between density per area and density per solid angle.
\clearpage
\section*{Verification case: probability densities carry a measure}
\noindent A uniform random point on a rectangular surface of area
\(6\,\mathrm m^2\), with width \(2\,\mathrm m\) and height
\(3\,\mathrm m\), has area-density
\(p_A=1/6\,\mathrm m^{-2}\). Parameterize the surface with
dimensionless \(s,t\in[0,1]\):
\[
x=2s\ \mathrm m,\qquad y=3t\ \mathrm m.
\]
Its change-of-variables Jacobian has area factor
\(\left|\frac{\partial(x,y)}{\partial(s,t)}\right|=6\,
\mathrm m^2\), so the unit-square parameter density is
\(p_{s,t}=p_A\,6\,\mathrm m^2=1\).
The different numbers express the same distribution relative to
different measures. Confusing a density with a probability can make
otherwise plausible sampling weights dimensionally incorrect.

For a small visible surface patch at distance \(r\) and incidence
angle \(\theta\), its solid angle is
\(d\omega=|\cos\theta|\,dA/r^2\), where the local geometry
defines the angle consistently. Therefore an area-density
\(p_A\) transforms to a solid-angle density
\[
p_\omega=\frac{p_A\,r^2}{|\cos\theta|}.
\]
This formula requires compatible domains, nonzero distance and a
well-defined non-grazing mapping. At grazing incidence the
Jacobian degenerates; one must not silently divide by zero.
A normalized density integrates to one only against its own
measure.

\paragraph{Reconstruction.}
Sample the unit square, map points onto the rectangle, and
verify the expected uniform density in each coordinate system.
Then evaluate a visible planar patch from a fixed viewpoint and
numerically check the area-to-solid-angle Jacobian.
Document failures at degenerate geometry and compare any estimator's
weight \emph{and its units} before moving to path tracing.
